\documentclass[11pt,reqno]{amsart}
\usepackage{amsmath,amssymb,graphicx,amscd}
\usepackage{mathrsfs}
\usepackage{enumerate}
% pour les S des spheres
\usepackage{mathrsfs}
\usepackage{dsfont}
\usepackage{stmaryrd}
\usepackage{wrapfig}
\usepackage{pstricks}

% ----------------------------------------------------------------
\vfuzz2pt % Don't report over-full v-boxes if over-edge is small
\hfuzz2pt % Don't report over-full h-boxes if over-edge is small

\newcommand{\prob}{{\mathbb{P}}}

\renewcommand{\S}{\mathscr{S}}  % complex unit sphere

\newcommand{\CC}{{\mathbb C}}
\newcommand{\HH}{{\mathbb H}}
\newcommand{\RR}{{\mathbb R}}
\newcommand{\QQ}{{\mathbb{Q}}}
\newcommand{\NN}{{\mathbb{N}}}
\newcommand{\ZZ}{{\mathbb{Z}}}
\renewcommand{\P}{{\mathbb{P}}}
\renewcommand{\span}{{\mathrm{span}}}
\newcommand{\sgn}{\operatorname{{sgn}}}
 \DeclareMathOperator{\Prob}{\mathbb{P}}   %probability
 \DeclareMathOperator{\E}{\mathbb{E}}      %expectation
 \DeclareMathOperator{\cov}{cov}             %logarithmic integral
\newcommand{\1}{\mathds{1}}

 \newcommand{\I}{1\!\!1}                   %Indicator function


 \newcommand{\rd}{{\mathrm{d}}}            %differential
 \newcommand{\ii}{{\mathrm{i}}}
  \newcommand{\Id}{{\mathrm{Id}}}
\newcommand{\law}{\overset{\mbox{\rm \scriptsize law}}{=}}
\newcommand{\convlaw}{\overset{\mbox{\rm \scriptsize law}}{\longrightarrow}}


\renewcommand{\Re}{{\mathfrak{Re}}}
\renewcommand{\Im}{{\mathfrak{Im}}}

% THEOREMS -------------------------------------------------------

\newtheorem{thm}{Theorem}[section]
\newtheorem{cor}[thm]{Corollary}
\newtheorem{lem}[thm]{Lemma}
\newtheorem{prop}[thm]{Proposition}
\theoremstyle{definition}
\newtheorem{defn}[thm]{Definition}
\theoremstyle{remark}
\newtheorem*{rem}{Remark}



\numberwithin{equation}{section}


\begin{document}

~\vspace{-3cm}



\begin{center}
{\bf Probability, homework 6, due October 21.}
\end{center}
\vspace{0.5cm}
\small







\noindent{\bf Exercise 1}. Let $(X_i)_{i\geq 1}$ be i.i.d. Gaussian with mean 1 and variance 3. Show that
$$
\lim_{n\to\infty}\frac{X_1+\dots+X_n}{X_1^2+\dots+X_n^2}=\frac{1}{4}\ {\rm a.s.}
$$

\noindent{\bf Exercise 2}. Let $f$ be a continuous function on $[0,1]$. Calculate the asymptotics, as $n\to\infty$, of 
$$
\int_{[0,1]^n}f\left(\frac{x_1+\dots+x_n}{n}\right)\rd x_1\dots\rd x_n.
$$



\noindent{\bf Exercise 3}. The goal of this exercise is to prove that any  function, continuous on an interval of $\RR$, 
can be approximated by polynomials, arbitrarily close for the $L^\infty$ norm (this is the Bernstein-Weierstrass theorem).  
Let $f$ be a continuous function on $[0,1]$. The $n$-th Bernstein polynomial is
$$
B_n(x)=\sum_{k=0}^n\dbinom{n}{k}x^k(1-x)^{n-k} f\left(\frac{k}{n}\right).
$$

a) Let $S_n(x)=B^{(n,x)}/n$, where $B^{(n,x)}$ is a binomial random variable with parameters $n$ and $x$:
$B^{(n,x)}=\sum_{\ell=1}^n X_i$ where the $X_i$'s are independent and $\Prob(X_i=1)=x$, $\Prob(X_i=0)=1-x$.
Prove that $B_n(x)=\E(f(S_n(x)))$.

b) Prove that $\|B_n-f\|_{L^\infty([0,1])}\to 0$ as $n\to\infty$.\\





\noindent{\bf Exercise 4}. Calculate
$
\lim_{n\to\infty} e^{-n}\sum_{k=0}^n\frac{n^k}{k!}.
$\\


\noindent{\bf Exercise 5.}
Let $(X_n)_{n\geq 0}$ be a sequence of i.i.d random variables, with uniform distribution on $[0,1]$. Let $Y_n=(X_n)^n$.
\begin{enumerate}
\item  Calculate the distribution of $Y_n$.
\item Show that $(Y_n)_{n\geq 0}$ converges to 0 in probability.
\item Show that $(Y_n)_{n\geq 0}$ converges in ${\rm L}^1$.
\item Show that almost surely  $(Y_n)_{n\geq 0}$ does not converge.
\end{enumerate}

\noindent{\bf Long problem}.
The goal is to prove the Erd{\H o}s-Kac theorem: if $w(m)$ denotes the number of distinct prime factors of $m$ and $k$ is a random variable uniformly distributed on $\llbracket 1,n\rrbracket$, then the following convergence in distribution holds:
$$
\frac{w(k)-\log\log n}{\sqrt{\log\log n}}\underset{n\to\infty}{\longrightarrow} \mathscr{N}(0,1).
$$
\begin{enumerate}[(i)]
\item Prove that if $(X_n)_{n\geq 1}$ converges in distribution to $\mathscr{N}(0,1)$ and $\sup_{n\geq 1}\E[X_n^{2k}]<\infty$ for any $k\in\mathbb{N}$, then $\lim_{n\to\infty}\E[X_n^k]=\E[\mathscr{N}(0,1)^k]$ for any $k\in\mathbb{N}$.
\item Prove that for any $x\in\mathbb{R}$ and $d\geq 1$ we have $$\left|e^{\ii x}-\sum_{\ell=0}^d\frac{(\ii x)^\ell}{\ell!}\right|\leq \frac{|x|^{d+1}}{(d+1)!}.$$
\item Assume that $\lim_{n\to\infty}\E[X_n^k]=\E[\mathscr{N}(0,1)^k]$ for any $k\in\mathbb{N}$.  Prove that $X_n$ converges in distribution to $X$. 
\item Let $w_y(m)$ be the number of prime factors of $m$ which are smaller than $y$.  Let $(B_p)_{p\ {\rm prime}}$ be independent random variables such that $\mathbb{P}(B_p=1)=1-\mathbb{P}(B_p=0)=\frac{1}{p}$,  $W_y=\sum_{p\leq y} B_p$, $\mu_y=\sum_{p\leq y}\frac{1}{p}$, $\sigma_y^2=\sum_{p\leq y}(\frac{1}{p}-\frac{1}{p^2})$. Prove that if $y=n^{{\rm o}(1)}$, then for any $d\in\mathbb{N}$ we have
$$
\E\left[\left(\frac{w_y(k)-\mu_y}{\sigma_y}\right)^d\right]-\E\left[\left(\frac{W_y-\mu_y}{\sigma_y}\right)^d\right]\to 0
$$ 
as $n\to\infty$.
\item Conclude.
\end{enumerate}


\end{document} 